\documentclass[10pt,a4paper]{amsart}
\usepackage[margin=19mm]{geometry}
\usepackage{amsmath,amssymb,mathtools}
\usepackage[scr=rsfs,cal=euler]{mathalfa}
\usepackage{xfrac}
\usepackage[unicode,hidelinks,bookmarks=false]{hyperref}
\newcommand{\ie}{i.e.\ }
\newcommand{\cf}{cf.\ }
\newcommand{\resp}{respectively\ }
\newcommand{\Subsection}{\subsection}
\providecommand{\nobreakdash}{\nobreak\hbox{-}}
\setlength{\emergencystretch}{2em}
\multlinegap0pt
\usepackage{tikz}
\usetikzlibrary{cd}
\usepackage{amssymb}
\usepackage[shortlabels]{enumitem}
\usepackage{tikz}
\usepackage{algorithm}\usepackage{algpseudocode}
\newtheorem{lemma}{Lemma}
\newtheorem{theorem}{Theorem}
\title{Rank Distribution and Dynamics of Gram Matrices from Binary m-Sequences with Applications to LCD Codes}
\author{Hengfeng Liu, Chunming Tang, Cuiling Fan, Zhengchun Zhou}
\date{Source: arXiv:2604.26387v1 (CC BY 4.0)}
\begin{document}
\maketitle

\noindent\emph{Source and licence.} Rank Distribution and Dynamics of Gram Matrices from Binary m-Sequences with Applications to LCD Codes. Version arXiv:2604.26387v1 (CC BY 4.0), distributed by arXiv under the Creative Commons Attribution 4.0 International licence, http://creativecommons.org/licenses/by/4.0/, as recorded in the arXiv licence metadata for this exact version and shown on the abstract page https://arxiv.org/abs/2604.26387v1. Full licence notice: \texttt{licenses/arxiv-2604.26387-CC-BY-4.0.txt}.\par\smallskip

\noindent\textit{Editorial scope.} Counted unit: the article's theorem thm:galois-descent (source0.tex lines 935--946). The article's complete original proof of it is delivered with it (source0.tex lines 948--1080, one \texttt{proof} environment of 133 lines). No mathematics is written, repaired or re-derived: the statement and the proof are the article's own lines, in the article's own notation, and no retained line is substituted or re-flowed.  Retained uncounted as the source-local context the proof needs: lines 916--922.  Nothing was omitted from any retained span, so the package carries no omission record.  No retained line contains a citation, so the wrapper carries no bibliography and no bibliography entry.  The retained lines contain the article's own diagram code, which the wrapper typesets with the standard tikz packages; no image file is delivered and the package declares no asset dependency.  Editorial formatting only, in addition: an AMS \texttt{amsart} preamble that restates the article's own declarations and macros used by the retained lines, namely the environments \texttt{lemma}, \texttt{theorem} on separate counters, together with the standard packages those lines need.  The retained environments print in this document as Lemma 1, Theorem 1 in order of appearance, whereas the article prints the same items with the numbering of its own full document.  The complete native source of the article as distributed in the arXiv e-print for this exact version is bundled unchanged as \texttt{sources/199391/source0.tex}.
\begin{lemma}[Dedekind's lemma]\label{lem:dedekind}
	Let $G$ be a group, $L$ be a field, and let $\chi_1, \dots, \chi_n : G \to L^{\times}$ be pairwise distinct $L$-characters of $G$. If $\lambda_1, \dots, \lambda_n \in L$ satisfy
	$$
	\sum_{i=1}^{n}\lambda_i \chi_i(g) = 0, \quad \forall  g \in G,
	$$
	then $\lambda_i =0, \, 1\le i\le n$.
\end{lemma}

\begin{theorem}\label{thm:galois-descent}
	Let $L/K$ be a finite Galois extension with Galois group $G = \mathrm{Gal}(L/K)$, and let
	$V$ be an $L$-vector space equipped with a semilinear representation $\rho : G\to \Gamma\rm{L}(V)$, $g\mapsto \rho_{g} $. Define a $K$-subspace of $V$:  $$V^G = \{v \in V \mid \rho_{\sigma}(v) = v, \, \forall \sigma \in G\}.$$ Then the
	map
	\[
	\mu \colon V^{G}\otimes_{K}L \longrightarrow V,
	\quad
	v\otimes \lambda \longmapsto \lambda v,
	\]
	is an isomorphism of $L$-vector spaces. In particular,
	$\dim_{K}V^{G}=\dim_{L}V$.
\end{theorem}

\begin{proof}
We first prove the injectivity. Suppose, for contradiction, that $\ker(\mu)\neq 0$, and pick a non-zero element
	\[
	\omega=\sum_{i=1}^{r}v_{i}\otimes \lambda_{i}\in\ker(\mu)
	\]
	with $r$ chosen as the smallest. Clearly $\lambda_{i}\neq 0$ for each $i$. Dividing $\omega$ by $\lambda_{1}$,
	we may further assume $\lambda_{1}=1$. The condition $\mu(\omega)=0$
	reads
	\begin{equation}\label{eq:inj-star}
		\sum_{i=1}^{r}\lambda_{i}v_{i}=0.
	\end{equation}
	Applying any $\sigma\in G$ to Eq. \eqref{eq:inj-star} and together with $\sigma(v_{i})=v_{i}$, we obtain
	\begin{equation}\label{eq:inj-starstar}
		\sum_{i=1}^{r}\sigma(\lambda_{i})\,v_{i}=0, \quad \forall\,\sigma\in G.
	\end{equation}
	Subtracting Eq. \eqref{eq:inj-star} from Eq. \eqref{eq:inj-starstar} yields
	\begin{equation}\label{eq:inj-dagger}
		\sum_{i=2}^{r}\bigl(\sigma(\lambda_{i})-\lambda_{i}\bigr)\,v_{i}=0,
		\quad \forall\,\sigma\in G.
	\end{equation}
	
	Suppose that some $\lambda_{i_{0}}$ with
	$i_{0}\geq 2$ does not lie in $K=L^{G}$. Then there exists
	$\sigma_{0}\in G$ with
	$\sigma_{0}(\lambda_{i_{0}})\neq \lambda_{i_{0}}$. Specializing Eq.
	\eqref{eq:inj-dagger} to $\sigma=\sigma_{0}$ gives that the element
	\[
	\omega':=\sum_{i=2}^{r}v_{i}\otimes
	\bigl(\sigma_{0}(\lambda_{i})-\lambda_{i}\bigr) \in V^{G}\otimes_{K}L
	\]
	satisfies $\mu(\omega')=0$ and has its $i_{0}$-th coefficient non-zero,
	so $\omega'\neq 0$; but this contradicts the
	minimality of $r$.
	
	Therefore $\lambda_{i}\in K$ for every $i$, together with Eq. \eqref{eq:inj-star}, we obtain
	\[
	\omega=\sum_{i=1}^{r}v_{i}\otimes \lambda_{i}
	=\sum_{i=1}^{r}(\lambda_{i}v_{i})\otimes 1
	=\Bigl(\sum_{i=1}^{r}\lambda_{i}v_{i}\Bigr)\otimes 1=0,
	\]


contradicting the assumption $\omega\neq 0$. Hence $\ker(\mu)=0$.
	
	
Next, we prove the surjectivity of $\mu$. Let $W\subseteq V$ be the image of $\mu$, which is the $L$-span of $V^{G}$. We prove that $W=V$.
	
Observe that $W$ is stable under the $G$-action on $V$, i.e., $\rho_{\sigma}(W)\subseteq W$. Indeed, for any
$\sigma\in G$, and any element $\sum_{j}\ell_{j}w_{j}$ of $W$ with $\ell_{j}\in L$ and $w_{j}\in V^{G}$, we have
\[
\rho_{\sigma}\Bigl(\sum_{j}\ell_{j}w_{j}\Bigr)
=\sum_{j}\rho_{\sigma}(\ell_{j}w_{j})
=\sum_{j}\sigma(\ell_{j})\,w_{j}\in W.
\]


Since $W$ is an $L$-subspace of $V$, the quotient $V/W$ is an
$L$-vector space. Since $W$ is a $G$-stable $L$-subspace of $V$, it is a subrepresentation
of $V$. The quotient $V/W$ therefore inherits the structure of a
semilinear $G$-representation, the quotient representation,
whose action $\bar\rho\colon G\to \Gamma L(V/W)$ is uniquely determined
by the commutativity of
\[
\begin{tikzcd}
	V \arrow[r, "\rho_{\sigma}"] \arrow[d, "\pi"'] & V \arrow[d, "\pi"] \\
	V/W \arrow[r, "\bar\rho_{\sigma}"'] & V/W
\end{tikzcd}
\]

for every $\sigma\in G$, where $\pi\colon V\to V/W$ is the canonical
projection. Explicitly, $\bar\rho_{\sigma}(\overline{v})
=\overline{\rho_{\sigma}(v)}$, and the semilinearity relation
$\bar\rho_{\sigma}(\ell\,\overline{v})=\sigma(\ell)\,
\bar\rho_{\sigma}(\overline{v})$ descends from the corresponding
relation on $V$.

Now suppose, for contradiction, that $V/W\neq 0$, and fix any
$u\in V\setminus W$, so that $\overline{u}\neq 0$ in $V/W$. For each
$\lambda\in L$, consider
\[
\widetilde{T}(\lambda):=\sum_{\sigma\in G}\rho_{\sigma}(\lambda u)
\in V.
\]
For every $\eta\in G$, we have
\[
\rho_{\eta}\bigl(\widetilde{T}(\lambda)\bigr)
=\sum_{\sigma\in G}\rho_{\eta\sigma}(\lambda u)
=\widetilde{T}(\lambda),
\]
so $\widetilde{T}(\lambda)\in V^{G}\subseteq W$. Passing to $V/W$ and
using semilinearity of $\bar\rho_{\sigma}$, the image
$\overline{\widetilde{T}(\lambda)}=0$ reads
\begin{equation}\label{eq:surj-star}
	\sum_{\sigma\in G}\sigma(\lambda)\,\bar\rho_{\sigma}(\overline{u})=0
	\quad\text{in } V/W, \quad \forall\,\lambda\in L.
\end{equation}

Fix an $L$-basis $\{\bar e_{k}\}_{k\in\mathcal K}$ of $V/W$, and
expand
\[
\bar\rho_{\sigma}(\overline{u})=\sum_{k\in\mathcal K}c_{\sigma,k}\,
\bar e_{k}, \quad c_{\sigma,k}\in L,
\]
where for each fixed $\sigma$ only finitely many $c_{\sigma,k}$ are
non-zero. Substituting into Eq. \eqref{eq:surj-star} and reading off the
coefficient of $\bar e_{k}$ gives, for every $k\in\mathcal K$ and
every $\lambda\in L$,
\begin{equation}\label{eq:surj-char}
	\sum_{\sigma\in G}c_{\sigma,k}\,\sigma(\lambda)=0.
\end{equation}

For each $\sigma\in G$ we regard the
restriction
\[
\chi_{\sigma}:=\sigma|_{L^{\times}}\colon L^{\times}\longrightarrow L^{\times}
\]
as an $L$-character of the group $L^{\times}$, then distinct $\sigma\in G$ yield distinct characters
$\chi_{\sigma}$. Restricting Eq. \eqref{eq:surj-char} to
$\lambda\in L^{\times}$ yields
\[
\sum_{\sigma\in G}c_{\sigma,k}\,\chi_{\sigma}(\lambda)=0,
\quad \forall\,\lambda\in L^{\times}.
\]
 Applying
Lemma~\ref{lem:dedekind} forces $c_{\sigma,k}=0$ for every
$\sigma\in G$ and every $k\in\mathcal K$. Hence
$\bar\rho_{\sigma}(\overline{u})=0$ in $V/W$ for every $\sigma\in G$. Taking $\sigma=\mathrm{id}_{L}$, so that
$\bar\rho_{\mathrm{id}_{L}}=\mathrm{id}_{V/W}$, yields
$\overline{u}=0$, contradicting the choice of $u\in V\setminus W$.
Hence $V=W$, and $\mu$ is surjective.
	

\end{proof}

\end{document}
